\chapter{De l'individu à la communauté patrimoniale pirate}

\begin{quote}
    \textit{L'autre inconvénient de ce type de bibliothèque est qu'elle permet, appelle, encourage la xérocivilisation. Cette civilisation de la photocopie, en dépit des grands avantages qu'elle procure, entraîne avec elle une série de dangers pour l'édition, même d'un point de vue légal. C'est avant tout la mort de la notion de droit d'auteur. [...] Sans compter qu'on peut emprunter le livre, le sortir et que certaines coopératives d'étudiants vous feront les photocopies sur des feuilles perforées prêtes à ranger dans des classeurs. Dans ces coopératives aussi on vous dit parfois qu'on ne peut pas photocopier un livre entier : j'ai eu ce problème avec certains de mes  étudiants. Ils me disent : \og Il nous faut trente copies de ce livre mais ils refusent \fg.[...] Je leur dis : \og Très bien, faites faire une photocopie, puis rapportez le livre à la bibliothèque ; après demandez vingt-neuf copies d'une photocopie. Il n'y a pas de droits sur une photocopie\fg.}\\

    Umberto Eco, \textit{De Bibliotheca}, traduit de l'italien par Eliane Deschamps-Pria, édition de l'Echoppe, 1986.
\end{quote}

Dans ce chapitre, nous nous concentrerons sur les origines du lien entre les individus, les pratiques du piratage, et les communautés patrimoniales pirates. Nous nous intéresserons d'abord à la naissance des carrières pirates, puis nous comparerons les différentes manières d'apprendre et de s'autonomiser chez les nouveaux pirates. Enfin, nous terminerons avec l'utilisation des compétences acquises par le piratage dans le monde professionnel et académique. Nous nous inscrirons ici dans l'héritage des travaux de Patrice Flichy, sociologue des sciences de l'information et de la communication ayant étudié la massification et l'utilisation d'internet chez les amateurs \footnote{Patrice Flichy, \textit{Le sacre de l'amateur, sociologie des passions ordinaire à l'ère numérique}, collection République des idées, éditions du Seuil, 2014}. La grande majorité de nos enquêtés sont des membres de la \textit{génération Z}\footnote{la génération Z comprend les personnes nées entre la fin des années 1990 et et le début des années 2010.} qui ont grandi avec le développement de l'informatique et de la télécommunication domestique grand public. L'étude de l'acquisition des compétences des pirates nous plongera également dans l'apparition massive de la sociabilisation en ligne.

\section{Ohé Moussaillons ! Dynamiques de découverte et de transmission chez les jeunes pirates}

Dans cette première partie, nous étudierons les moments de découverte du piratage chez nos enquêté.es, nous nous intéresserons spécifiquement à la dimension sociale et émotionnelle des premiers pas dans l'univers du piratage.

\subsection*{Définir la culture pirate}

Pour qu'un individu commence à pirater, il doit alors passer par une étape de découverte et d'acceptation d'un nouvel usage, et dans le cas du piratage d'une nouvelle culture. Même très jeunes, nos enquêté·es connaissent les limites légales de la copie, et savent que l'utilisation qu'ils en font est illégale. Ils se retrouvent donc face à un dilemme moral : respecter la loi ou choisir de l'enfreindre consciemment, et nous allons essayer de comprendre les éléments qui fondent leurs choix. Nous pouvons résumer la culture pirate en quelques clés \footnote{\cite{de_kosnik_rogue_2016} page 2}.

\begin{enumerate}
    \item La culture pirate est l'assomption que tout est possible sur le réseau, de la citation, la transformation et le partage de biens culturels, sans citer et sans en demander l'autorisation à l'auteur. Ce cadre de pensée est également la base d'objets culturels comme les \textit{les memes} dont le principe est de tordre et réinventer un bout d'objet culturel sorti de son contexte à l'infini, les \textit{fanfictions} qui prennent un univers existant pour le transformer selon les envies des auteurices, ou les \textit{mods} dont le but est d'ajouter des fonctionnalités à un jeu, quitte à le transformer radicalement.
    \item Le piratage n'est pas assimilable à du vol, ni à un détournement de propriété, mais bien à la production non autorisée d'une copie numérique d'un contenu culturel. De ce fait, il existe une différence entre le cadre légal et l'éthique, et l'un ne définit pas l'autre\footnote{Richard Stalman,
\textit{Annexe 4 : Quelques expressions trompeuses qu'il vaut mieux éviter d'employer}, dans O. Blondeau, Libres enfants du savoir numérique : Une anthologie du “Libre” (p. 359-364). Éditions de l'Éclat, définition de vol.}
\end{enumerate}

La découverte des pratiques de piratage de biens culturels est un parcours relativement homogène chez mes enquêtés. Sauf quelques rares exceptions, les individus découvrent le piratage avec leurs premiers équipements numériques (premier téléphone, premier mp3, ordinateur familial, premier pc). Cette découverte passe par un rapport entre adolescents qui ont les mêmes équipements, ou par une personne extérieure (parents, amis des parents, grands frères et sœurs) que l'on peut surnommer comme étant des \textit{figures tutélaire}. L'informatique devient alors l'interface et le principal médium de l'intérêt pour la musique, le cinéma, les séries, les jeux vidéo et les mangas.

\subsection*{La découverte par les pairs}

La voie majoritaire de découverte du piratage se passe dans le cadre amical durant l'enseignement secondaire. En discutant avec ses ami.es de contenus culturels, principalement de musique, l'individu prend conscience de l'existence d'outils et de communautés dédiés à l'acquisition pirate de ses centres d'intérêts. Ce chemin d'apprentissage privilégie l’acquisition de compétences de surface, permettant l'acquisition de contenus en passant par des services tiers. Youtubemp3, un site de conversion de vidéo YouTube en fichiers mp3 aujourd'hui disparu, mais dont des milliers de clones ont pris la relève, a été pour la majorité de mes enquêtés la première étape de la découverte du piratage. La principale méthode de diffusion de ce site est alors la discussion entre ami.es durant la récréation du collège. La découverte de l'outil creuse ensuite la curiosité au fur et à mesure que les goûts culturels de l'individu s'autonomisent et que l'informatique se démystifie petit à petit.

Hilaire et Élise sont deux profils ayant commencé à pirater de cette manière. Hilaire s'est formé au piratage avec ses groupes de fans et de partage de dramas coréens traduits. Le cœur de l'activité de pirate d’Élise réside dans le téléchargement de livres, essentiellement des classiques, des livres qu'elle possède déjà, ou des livres de romance qu'elle n'oserait pas acheter en magasin.

\begin{quote}
H - \textit{Ma grand-mère m'avait offert un ordi portable qui coûtait très cher à l'époque.
 [...] Au début, je pense que ça a commencé par la musique, je pense que c'est le tout premier truc. C'est genre les sites où tu mets le lien ... Youtubemp3 un truc comme ça. [...] J'ai encore tous mes fichiers, donc ça, c'est un peu marrant, je pourrais aller voir les tout premiers trucs que j'ai enregistrés.}
\end{quote}

\begin{quote}
    E - \textit{Déjà, quand j'ai commencé à pirater, j'avais un ami. Au départ ce n'était pas vraiment moi qui le faisais. J'ai le père d'une amie proche qui, lui, a téléchargé énormément de films et tout, et il me les passait sur clé USB. }
\end{quote}

La découverte du piratage entre pairs favorise une découverte graduelle et relativement lente de solutions peu techniques. Ici, la conceptualisation de l'apprentissage du piratage ressemble à une découverte géographique plutôt qu'à des compétences informatiques. Les individus se partagent des \og bonnes adresses \fg et nous pouvons résumer leur horizon mental à une carte s'enrichissant avec le temps de lieux (des sites internet). Cependant, cette capacité de recherche de solutions fait que ces gens ont en général un niveau satisfaisant de maîtrise de la technologie \footnote{Unesco Regional Bureau for Education in Asian and Pacific \textit{Diverse approaches to developing and implementing competency-based ICT training for teachers: a case study}, TH/EISD/16/019-500, 2016}, même si cela n'est que le médium entre eux et leurs passions. Ces gens entrent rarement en contact avec des communautés patrimoniales pirates, sauf pour le cas particulier du piratage académique dont nous allons reparler un peu plus loin.\\

%A compléter et enrichir
Dans les archives des skyblogs de la BnF, nous trouvons plusieurs blogs dont le but est de partager des mangas ou des films que l'on a appréciés, tout en partageant également un lieu de téléchargement ou de streaming pour ses ami.es et followers. Par exemple, le blog de Spymon : \textit{telechargement09} \footnote{voir Annexes Skyblogs} qui se concentre sur le partage de sa passion pour la musique et l'esthétique du genre métal/gothique. Nous pouvons remarquer que les personnes téléchargeant ou partageant de la musique au format mp3 sur internet ont souvent le réflexe de ne pas assumer cette action comme du piratage. C'est également produit dans un de nos entretiens où Élise, cherchant dans sa mémoire l'instant où elle a commencé à pirater, s'est rendu compte au fur et à mesure des exemples que sa première rencontre avec la piraterie était le téléchargement de musique. 

\subsection*{La découverte par une figure tutélaire}

La deuxième manière dont on peut découvrir le piratage est en passant par une \textit{figure tutélaire}, une personne déjà très à l'aise avec l'informatique, la technologie et le piratage, qui va prendre le rôle de messager vertical entre les communautés patrimoniales pirates et l'individu. Cette manière de découvrir le piratage peut imprimer deux manières de faire chez les personnes qui commencent à construire une habitude de piratage. La première, sans doute la plus commune, est un apprentissage plus complet et une montée en compétences très rapide de ces individus. On y retrouve des pratiques comme le partage en pair à pair, le torrent, le partage d'applications modifiées, la mise en place de serveurs, de seedbox, ou de VPN comme Hamachi pour faciliter le partage ou le lien avec ses ami.es. L'exemple type est celui d'Adrien, qui a été introduit aux pratiques de piratage très jeune grâce à son père qui lui apprend et lui finance des moyens de pirater dans une sécurité relative.

\begin{quote}
    A - \textit{Donc depuis tout petit, moi je vois des ordinateurs à la maison, [...] avec ma mère j'ai découvert l'informatique et la bureautique, et avec mon père plus tard j'ai découvert le reste. Donc beaucoup de jeux vidéos, et aussi de la consommation de contenu, illégal particulièrement. C'est mon père qui m'a fait découvrir le monde merveilleux de la R4\footnote{La R4, pour \textit{Revolution for DS} est un \textit{linker}, une cartouche permettant d'ajouter énormément de fonctionnalités à une console, en particulier tricher, ou faire fonctionner des ROMs de jeux sans avoir de cartouche officielle. S'il n'existe pas de loi encadrant l'existence des linkers et leur utilisation, des revendeurs ont été condamnés en France, voir l'article de Julien Lausson, \textit{La vente de linkers pour Nintendo DS jugée illégale en France}, Numérama, 04/10/2011} La vente de linkers pour Nintendo DS jugée illégale en France  à l'époque, et surtout le site EmulIsland, je crois que ça s'appelait, pour les jeux Wii, je m'en souviens encore. Et donc ça c'était avec mon père que j'ai découvert ça.}
\end{quote}

Une autre évolution possible de la découverte par une figure tutélaire peut être le fait de continuer de se reposer sur l'expérience de cette personne. Ceci semble arriver chez des gens qui n'ont pas l'envie, le temps, ou la curiosité d'approfondir leur rapport à l'informatique. C'est par exemple le cas de Zakary avec son frère. Les deux ont toujours été très proches malgré la différence d'âge et l'un des principaux vecteurs de cette complicité fraternelle est l'amour de la  musique et du cinéma d'animation japonais.

\begin{quote}
    Z - \textit{J'ai un grand frère [...] et c'est cette personne-là qui m'a introduit au piratage. Quand j'étais très jeune, un peu avant le collège, il m'a un peu expliqué sur l'ordinateur familial comment récupérer les vidéos YouTube que j'aimais bien, comment mettre des musiques sur mon téléphone portable pour que je puisse les écouter quand je voulais. [...] Mon frère avait certains intérêts que moi, je n'avais pas forcément. Et donc, quand je voulais ouvrir une porte vers ces mondes-là, je n'avais pas forcément les codes. Je savais que j'avais un peu son héritage qui pouvait être là. J'ai même bien parlé avec lui pour qu'il me donne quelques clés.}
\end{quote}

L'un des points culminants de cette relation trouve son incarnation dans un objet, un iPod rempli de ses musiques piratées, que le frère de Zakary lui offre. Ce qui semble expliquer que Zakary n'ait pas eu le même parcours qu'Adrien est qu'il est plus intéressé par les goûts musicaux et cinématographiques de son frère que par les méthodes et la \textit{sociabilisation pirate}.
    \begin{quote}
        Z - \textit{Cela a définit mes goûts musicaux pour le restant de mes jours.}
    \end{quote}

\subsection*{La découverte en autonomie}

Une troisième manière de découvrir le piratage nous est indiquée dans le vécu de Théo, qui a grandi dans un lieu isolé en pleine campagne, avec l'accès à un ordinateur dès la classe de sixième. Avec des parents qui utilisaient un minitel pour des tâches administratives, la famille de Théo est une famille qui connaît et utilise des services informatiques, sans pour autant avoir une figure tutélaire technophile. La spécificité du parcours de Théo est l'autoformation, conséquence immédiate de sa passion naissante pour l'informatique qu'il partage avec son frère, qui lui offre la connexion à un monde culturel qui lui manque.

\begin{quote}
    T - \textit{J'étais assez isolé en soi des grandes villes et tout [...] j'ai assez rapidement eu accès à un ordinateur, comment utiliser un ordinateur, comment utiliser internet, des cet âge là. Par contre mon téléphone portable je ne l'ai eu qu'à partir de la seconde. Je n'avais pas le droit avant, ça peut paraître un peu bizarre d'ailleurs maintenant sur l'évolution d'internet. [...] J'ai un peu trouvé ma passion là dessus sur l'informatique de manière générale ...}\\

    \textit{Mon utilisation au tout début c'était au lycée où je commençais à vouloir trouver des films. Je reviens au fait que je vivais en campagne mais en fait j'avais pas mon autonomie, et je pouvais pas aller au lycée, je pouvais pas aller au cinéma à pied ou en bus ou quoi. Donc ça dépendait de mes parents, mes parents n'étaient pas très intéressés par le cinéma, moi j'adorais ça. Du coup le but c'était d'aller chercher ce que je pouvais sur internet et donc assez rapidement je suis tombé sur des méthodes de téléchargement, comment essayer de m'en sortir tout seul ...}
\end{quote}

C'est cette combinaison de l'informatique comme un médium pour une passion pour le cinéma et une passion pour l'informatique qui permet à Théo de briser la barrière entre son isolement et les communautés pirates. \\


\section{Grandir en piratant et évolution des pratiques au fil de la vie}

\begin{quote}
    \textit{La démocratisation des compétences repose d'abord sur l'accroissement du niveau moyen de connaissances (dû notamment à l'allongement de la scolarité) et sur la possibilité offerte par Internet de faire circuler les savoirs, de livrer son opinion à un public plus vaste. L'amateur qui apparaît aujourd'hui à la faveur des techniques numériques y ajoute la volonté d'acquérir et d'améliorer des compétences dans tel ou tel domaine.}\\

Patrice Flichy, \textit{Le sacre de l'amateur, sociologie des passions ordinaires à l'air numérique}, 2014
\end{quote}

L'un des éléments les plus intéressants dans l'étude des carrières pirates est l'évolution de la pratique au cours de la vie de nos enquêté.es. En devenant de jeunes adultes, les individus se découvrent de nouvelles passions, et, entrant dans les études supérieures (tous mes enquêté.es ont fait des études supérieures jusqu'à minimum bac+3), iels réutilisent ce qu'ils faisaient pour les films et les livres dans le cadre académique. Nous nous concentrerons dans cette partie sur l'évolution et l'adaptation des individus qui ont appris à pirater durant leur adolescence et qui vont continuer leur carrière pirate dans leurs études, leur vie professionnelle ou personnelle. Nous déshabillerons leurs routines, leurs justifications, et les évolutions qu'iels ont dû mettre en place pour s'adapter à un contexte qui change. % reformuler la dernière phrase de merde

\subsection*{Pérennisation, transformation ou abandon des habitudes pirates des jeunes adultes}

La vie de jeune adulte semble produire deux conséquences sur les habitudes de piratage de nos enquêté.es. Certains, comme Théo et Adrien, semblent creuser et intensifier leurs habitudes de piratage ; d'autres, comme Hilaire et Élise, semblent mettre en pause pendant un temps leurs carrières de pirates. Ces moments de pause semblent liés à l'émergence de nouveaux types de produits culturels comme la vidéo à la demande et les abonnements de streaming dont le représentant principal est Netflix. La praticité et le coût très faible ont permis à certaines entreprises, comme Netflix, Spotify ou Steam, de remplacer les habitudes des pirates par des alternatives légales. Le changement de rythme de vie, de nouveaux temps de transport, le fait d'avoir plus de moyens à mettre dans sa consommation culturelle, sont diverses raisons d'abandon des carrières de pirates des jeunes adultes. Le téléchargement demande du temps et une organisation qui ne peuvent pas être compatibles avec la vie d'une personne s'insérant professionnellement ou qui habite loin.  

\begin{quote}
    H - \textit{J'ai un abonnement Netflix, et je pense que ouais du coup à cette période j'arrête totalement de pirater films, séries, j'ai un peu continué pour de la musique durant ma licence, mais ça n'a pas duré très longtemps parce qu'après j'ai eu un compte Spotify. Et c'était quand même vachement pratique de ne pas avoir à aller chercher tous les trucs que tu as envie d'écouter, de les télécharger et de mettre sur ton téléphone. À la limite je pense que j'ai dû aller quelques fois sur des sites pour lire des scans de mangas. Mais ça n'a pas été vraiment très très dur, mais ça m'est pas arrivé beaucoup. }
\end{quote}

L'évolution d'une pratique de piratage peut signifier l'arrêt ou le ralentissement d'une partie de l'activité d'un individu, sans pour autant signifier l'essoufflement total de la carrière de pirate. Par exemple, Élise pirate énormément de livres, et continue de pirater des livres, mais elle a arrêté de télécharger des livres audio, auxquels elle accède maintenant par un service légal qu'elle paie. 

\begin{quote}
    E - \textit{Moi, c'est venu vraiment après, quand j'ai découvert le monde merveilleux du piratage de livres. Même des livres audio, des fois, j'en prenais aussi. Ça, d'ailleurs, c'est un autre moyen puisque les livres audio, avant, je les piratais. Mais maintenant il y en a plein sur Spotify, par exemple. Donc,il y a même des choses que je pensais illégales. Maintenant je les fais ... j'ai un abonnement Spotify chaque mois. Mais je peux lire, je peux écouter les Hauts de Hurle-Vents, je sais pas, enfin gratuit quoi, enfin \og gratuit \fg entre guillemets. Mais ça c'est venu récemment, je crois.}
\end{quote}

Le sujet animant toute la communauté travaillant sur le piratage est l'impact d'offres légales jugées supérieures à des expériences pirates \footnote{Anthony Galluzo à toute une sous partie sur cette question de l'impacte des offres légales de streaming sur les comportements illégaux.}. Cependant, cette question semble beaucoup plus sensible au niveau de revenu des individus que l'existence intrinsèque d'offres légales. Par exemple, avant l'apparition de Netflix, il existait durant les années 1990-2000 des magasins de prêt-vidéos permettant un accès légal à des films pour une fraction du prix de l'achat d'un DVD. L'apparition des plateformes de streaming n'a pas révolutionné le cœur de la question du piratage, mais elles ont cependant modifié les échelles auxquelles ces questions se posent. \\

Hilaire a un parcours particulier, après avoir presque totalement arrêté toute forme de piratage durant ses études supérieures et le début de sa vie professionnelle, il recommence à pirater durant sa reprise d'études d'un master recherche en sociologie. Sa reprise du piratage est ici également directement liée à l'impact de sa reprise d'études sur ses revenus. 

\begin{quote}
    H - \textit{Bha depuis ce temps bah du coup bon ça je te l'apprends pas mais au fil des années les plateformes Netflix, Spotify augmentent leur prix. Donc à un moment où moi ça m'a saoulé, du coup j'ai juste arrêté d'utiliser Netflix et Spotify? D'abord Netflix parce qu'en vrai c'était plus facile d'accéder à du streaming de séries et de films que je trouvais, et du coup la transition était assez facile, de revenir au piratage notamment par du torrent.[...] D'abord je suis passé sur Deezer parce que c'était ... enfin c'est une entreprise française et ils rémunéraient un peu mieux les artistes. Du coup j'ai d'abord fait ce choix en mode ... enfin un choix éthique quoi. Et après j'ai arrêté parce que je me suis rendu compte qu'en fait j'utilisais pas tant que ça Deezer. [...] Et récemment j'ai un ami qui m'a envoyé le lien pour télécharger une APK de Spotify et avoir Spotify sans pub, sans payer quoi. Et du coup j'ai un peu utilisé mais là ça commence à buger du coup il faut que je le retélécharge dans d'autres version pour voir si ça remarche. Donc je n'ai pas encore trouvé de solution pour la musique.}
\end{quote}

Nous reviendrons sur le point de jonction entre le vécu économique des individus et les arguments éthiques et politiques qu'ils déploient dans le chapitre suivant. Cependant, il est intéressant d'évaluer l'argument que nous avions déjà rencontré dans notre introduction, celui des pirates comme une manne de consommateurs potentiels. S'il n'est pas faux que des services comme Netflix, Spotify ou Steam aient réussi à générer un marché qui a converti une partie des populations pirates en clients, et que de nombreux spécialistes prédisaient la mort de la pratique du piratage ; la multiplication des plateformes et donc des abonnements, l'exclusivité du contenu, l'accès direct à la vidéo à la demande a redonné un souffle à la pratique. Les manières de pirater ont évolué en tenant compte de ces nouveaux acteurs et de l'apparition de l'économie de la publicité. En tapant le mot \og Torrent \fg dans l'outil Google Trends, nous possédons un témoin relativement fiable de l'intérêt du public pour le piratage utilisant des liens magnétiques depuis 2004 à nos jours \footnote{voir annexes Résultats Google Trends pour le mot Torrent}. Nous comprenons que la popularité du torrent s'effondre à partir de la deuxième moitié des années 2010 ; cependant, cela n'indique pas que les individus piratent moins mais que les méthodes de piratage évoluent. Une nouvelle manière de consommer du contenu pirate est le streaming, à la manière de Netflix ou Spotify. Les hébergeurs de contenus piratés ont adapté dans la grande majorité leur fonctionnement pour suivre les habitudes des consommateurs, abaisser les barrières à l'entrée (plus besoin de maîtriser le torrenting, plus besoin de télécharger un fichier), et rassurer les utilisateurs effrayés par l'association entre protocole Torrent et la surveillance d'Hadopi. Si Netflix et Spotify étaient des services révolutionnaires et peu onéreux pendant une époque, l'infrastructure économique et la superstructure culturelle de la société ont transformé cette pratique en une norme, et le fait social pirate, comme toutes les pratiques de déviance, a muté pour s'y insérer. 

\subsection*{La piraterie académique}

Dans les différentes manières dont les individus réutilisent, adaptent ou abandonnent leurs pratiques pirates, le piratage académique trouve une place particulière qui symbolise souvent le passage de l'adolescence (pirater pour de la consommation culturelle ludique) à la vie adulte et professionnelle (pirater pour accéder à de la connaissance et des compétences).

\begin{quote}
    \textit{A - Et donc ouais la prépa. Ça c'est un peu calmé. Je fais mes deux ans de classe préparatoire. Il y a quand même du contenu que j'ai téléchargé, qui est surtout en fait marginal, on avait la chance d'avoir une bibliothèque qui était bien fournie, donc on avait des manuels, mais il y avait certains manuels qu'on ne pouvait pas avoir autrement qu'en téléchargeant. Donc là, ça a été mon premier contact avec le téléchargement académique . Même si pas encore, parce que je ne savais pas tout à fait ce que c'était encore les articles de recherche vraiment, mais tout ce qui va être manuels d'histoire, manuel de maths, des trucs comme ça. Je ne pouvais pas lâcher 50 balles dans un manuel, et tu avais trois manuels par matière, ce n'était pas envisageable. Donc là, j'ai commencé à trouver des trucs sur internet. C'était bien.}
\end{quote}

Si le piratage académique accompagne le passage au monde adulte et professionnel, il accompagne donc également l'éveil des sensibilités politiques et des émotions patrimoniales. Les enquêtés ayant mentionné un rapport politique, une vision du monde ou une considération patrimoniale liée à la pratique du piratage l'ont intégré à leurs récits au moment de l'acquisition d'autonomie et de maturité des études supérieures et de l'entrée dans le monde professionnel. 

\begin{quote}
    T- \textit{C'est pas anodin effectivement [...] J'ai l'approche très scientifique de cette chose-là. C'est que pour moi, la science, en tout cas, doit être entièrement ouverte. J'ai beaucoup de mal avec le fait que les labos s'enferment sur des articles qui ne sont publiés qu'entre eux, que certains industriels ferment certains verrous sur l'accès à la science et à des technologies de pointes. Moi je suis pour, effectivement le partage complet, peu et simple, en fait, en science. Et du coup ça transpire un peu aussi sur ma conception du partage culturel. Je suis un gros soutien pour le fait que l'accès à la culture, en particulier pour les jeunes, doit être le plus simple possible. Donc ... en particulier sans frais, parce que pour les jeunes, forcément, le premier frein, ça doit être le budget. Donc le fait d'encourager les jeunes, que ce soit aller au musée, que ce soit aller au cinéma, que ce soit aller à une bibliothèque ou quoi, pour moi, tout doit être le plus simple possible.}
\end{quote}

Cette citation de Théo, après une demi-heure d'entretien, résonne avec des discours qui nous sont familiers, nous pouvons citer ici la page \textit{About} du site de médiation artistique pirate Ubuweb \footnote{Voir l'annexe Message d'introduction à Ubuweb}, ou avec le \textit{Manifeste pour une guérilla pour le libre accès} d'Aaron Swartz, ou du slogan du site Sci-Hub \footnote{« To remove all barriers in the way of science » (« Éliminer tout obstacle sur la voie de la science »)}. La plupart de nos enquêté·es ne se définissent pas comme des personnes particulièrement politisées. Si iels ont des idées et des sensibilités politiques, iels ne sont ni encarté·es, ni militant·es dans des organisations politiques. Cependant, la pratique du piratage amène un discours qui utilise un vocabulaire et des idées libertaires. Si ces personnes n'ont pas le temps et les moyens d'acquérir les compétences pour aligner leurs logiques et leurs actions, sociabiliser comme \og pirate \fg ou amateur, ou n'importe quel autre terme demande de rentrer en contact avec des communautés patrimoniales pirates, et donc de comprendre et de faire sienne une partie de la culture pirate et de ses influences comme la culture du libre et les idéologies libertaires. 

\begin{quote}
    H - \textit{Je ne me revendique pas de ça mais parce que le sujet ne me vient pas quoi. Enfin c'est pas trop quelque chose dont les gens autour de moi parlent. [...] Mais ouai je suis pas sûre que ce soit vraiment une identité même si en fait ... Ça fait depuis ... C'est en licence que j'ai découvert la culture du libre. Et j'avais eu des cours sur ça en fait par un intervenant qui travaillait dans un Fab Lab. Et je pense que vraiment ce cours a vraiment changé toute ma perspective de voir le monde. J'étais en mode ah mais oui c'est génial en fait. Genre tu peux hacker des machines, mais tu peux aussi hacker des vêtements, des trucs, enfin tu peux tout faire toi-même. [...] Ouai il nous a grave initié à la culture du libre. Wikipédia, Open Office, tout ça [...] Du coup c'est à cette époque que j'ai découvert un peu que c'était un mouvement de pirater des trucs, [...] j'avais des pratiques de piratage pas je savais pas que ça pouvait s'inscrire dans un mouvement, et dans des revendications politiques.}
\end{quote} 


\subsection*{De pirates à corsaires :  Quand le piratage s'invite dans le vie de l'entreprise}

Précédemment, nous avons exploré l'évolution des compétences de piratage dans le domaine académique et des études supérieures, mais cette mécanique s’étend au-delà des études. Pour certains, le piratage devient une facette de leur insertion professionnelle, où des entreprises recherchent et les rémunèrent pour leurs compétences. Certains milieux professionnels et secteurs économiques jouent avec les limites dans l'utilisation de certaines utilisations de contenu licencié. En discutant avec le comité de sélection \textit{L'Association des Archivistes Français}, une membre du jury me confie que le sujet de ce mémoire résonne particulièrement avec son vécu, les publications des normes AFNOR sont trop onéreuses pour son service d'archives, et son équipe a décidé de pirater ces documents pour pouvoir exercer correctement son travail. Cette même archiviste me parlera ensuite de l'affaire Filaé \footnote{Pierre Alexandre Compte, \textit{Archives : « L’affaire Filae », du besoin de repenser la politique de diffusion des données culturelles} dans la Gazette des communes, 21/12/2016}, anciennement Généalogie.com, une entreprise de généalogie ayant réalisé du profit sur la mise à disposition de documents numérisés et partagés gratuitement par les départements. J'aimerais introduire ici le mot \textit{bureautisation}, qui décrit à la fois l'économie de service tertiarisé et l'agencement physique des moyens de production au sein de l'espace du bureau, mais aussi de l'interface bureautique et donc des logiciels qui se font l'interface entre un.e travailleureuse et sa mission. La bureautisation amène des individus à privilégier l'alternative pirate : pour ne pas avoir à couper dans leurs budgets, pour ne pas avoir à faire une demande administrative chronophage à la hiérarchie, pour se tenir à jour des pratiques professionnelles de leurs branches.
Dans les milieux professionnels demandant un va-et-vient constant avec la recherche (milieu académique, professorat, ingénierie, médecine), le piratage devient le seul moyen de pouvoir accéder aux ressources pour que les étudiants, les apprentis et les professionnels puissent mettre à jour et progresser dans leurs compétences.

\begin{quote}
    Sylvie\footnote{Sylvie ne fait pas partie de nos enquêtés pour la simple et bonne raison qu'elle ne pirate pas, je la cite simplement suite à une discussion intéressante sur la place du piratage dans le monde des professions médicales} - \textit{Je donne des cours de français médical à des jeunes médecins immigrés, et je suis impressioné. Au début de l'année ils m'ont tous demandé le titre du manuel que j'utilise, et ils l'ont tous piraté, téléchargé, et ils se sont refilés les bonnes combines, c'est très impressionnant.}
\end{quote}

Le monde de l'intelligence artificielle, spécifiquement les grands modèles de langage (LLM) captent l'attention, les investissements et la couverture médiatique. Ces modèles demandant énormément de données en entrée, c'est logiquement que plusieurs entreprises ont été prises en train d'avoir recours au piratage de contenus culturels et scientifiques. L'exemple le plus récent est celui de la corporation de Mark Zuckerberg, Meta, un des acteurs les plus importants dans l’entraînement et la mise à disposition de LLM au grand public, qui a vu ses mails internes divulgués en démontrant un usage systématique de contenu licencié \footnote{Ashley Belanger, \textit{“Torrenting from a corporate laptop doesn’t feel right”: Meta emails unsealed}, sur le site arstechnica.com, 06/02/2025}. Meta s'est contenté de télécharger plus de 70 téraoctets de livres et papiers numérisés sur la plus grande bibliothèque pirate existante, Libgen. Les exemples d'acteurs privés épinglés pour du piratage culturel, ou du scraping de données personnelles dans le but de former des bases de données d’entraînement d'intelligence artificielle demanderont un mémoire à part entière pour être traités. 

\begin{quote}
    T - \textit{C'est compliqué à dire mais je l'ai même fait à la demande de mon employer assez étrangement, mais en fait c'est quelque chose qui est aussi répandu dans le milieu scientifique en particulier pour ce qui est de l'intelligence artificielle, le cœur du sujet, la base du travail dans l'intelligence artificielle c'est les bases de données. Si tu n'as pas de données tu ne peux pas entraîner de modèle donc tu ne peux rien faire et le problème c'est que les bases de données il faut les constituer. Il faut avoir des données propres, des données qui font du sens et du coup pour ça il faut avoir des données qui viennent de quelque part. Si ton entreprise n'est pas capable de générer elle-même les données qui l’intéresse il faut bien qu'elle aille les chercher ailleurs pour entraîner ses modèles et donc du coup très souvent tu vas avoir recours à du vol de données pur et simple. [...] Certains scientifiques et certaines entreprises ne se gênent pas pour aller voler ces jeux de données s'ils le peuvent [...] c'est que pour mon stage de fin d'études j'étais employé au sein de l'entreprise qui m'ont demandé expressément d'utiliser mon compte étudiant pour faire des demandes à des chercheurs qui sinon auraient fait payer les entreprises pour accéder à des données. J'ai eu aussi le cas de simplement aller essayer de piller des jeux de données sur des sites par forcément recommandables mais qui avait réussi à partager les jeux de données qui étaient utilisés pour différents articles.}
\end{quote}

\section{Conclusion}

La découverte du piratage se fait principalement par trois voies : la socialisation par les pairs, l'influence d'une figure tutélaire, et la curiosité personnelle. Chaque expérience apporte une perspective unique et offre des capacités d'apprentissage différentes aux individus. Il est important de préciser que ces différentes catégories s’interpénètrent et se superposent selon les situations. Ces interactions leur permettent d'acquérir des outils essentiels pour leur parcours académique ou professionnel, soulignant l'importance de la socialisation et de l'apprentissage continu dans le piratage. Comme le souligne Anthony Galluzo, la pratique du piratage est un ratio entre la capacité d'apprentissage, la mise en place de routines d'accès au contenu piraté, et l'entretien de ces compétences et routines. Les individus peuvent alors continuer ou reprendre leurs habitudes pirates et adapter les compétences qu’ils ont développées pour avoir accès à des ressources culturelles ou logicielles. Ce transfert de compétence se met en place durant les études supérieures et le besoin d’avoir accès à de la littérature scientifique, et peut continuer dans le monde professionnel en fonction de si la personne a besoin de continuer d’avoir accès à de la littérature scientifique, de la documentation technique, ou des outils logiciels et bases de données.

 
%III/ Dépôts numériques, hybride entre un lieu de
%conservation et de partage.
%• Le dépôt pirate, entre archive et bibliothèques
%– La fonction de conservation des dépôts numériques.
%– Le partage et la diffusion des ressources piratées.
%• Les communautés autour des dépôts numériques
%– Les contributeurs et les utilisateurs des dépôts.
%¨ Les dynamiques sociales et les interactions au sein des communautés.
%2
%– Les contributeurs et les utilisateurs des dépôts.
